%==============================================================================
\section{Evaluation Methodology}
\label{sec:evaluation}
%==============================================================================

We design an evaluation framework that combines automated structural metrics with human expert assessment, supported by an observability infrastructure that captures every agent interaction.

%----------------------------------------------------------------------
\subsection{Scoring Dimensions}
%----------------------------------------------------------------------

We define 16 scoring dimensions across four capabilities, each using a 5-point categorical scale $\mathcal{C} = \{\text{excellent}, \text{good}, \text{acceptable}, \text{poor}, \text{unacceptable}\}$ mapped to numeric values $\{1.0, 0.75, 0.5, 0.25, 0.0\}$:

\textbf{Threat Statement Quality} (5 dimensions): overall quality, relevance to application context, completeness of threat coverage, technical accuracy, and hallucination detection.

\textbf{Attack Tree Quality} (6 dimensions): overall quality, structural quality (depth, branching, organization), technical realism of attack techniques, logical progression of attack paths, completeness of attack vector coverage, and actionability for defenders.

\textbf{TTP Mapping Quality} (1 dimension): correctness of MITRE ATT\&CK technique assignment, using a specialized scale replacing ``unacceptable'' with ``no\_mapping'' for steps where no technique applies.

\textbf{Mitigation Quality} (4 dimensions): actionability (concrete implementable steps), specificity to the identified threat and technology stack, coverage of identified attack paths, and technical accuracy of recommended controls.

%----------------------------------------------------------------------
\subsection{Observability Infrastructure}
%----------------------------------------------------------------------

\threatforest{} produces two complementary trace types:

\textbf{OTEL traces} (automatic). The Strands SDK~\cite{strands2025} emits OpenTelemetry spans for every LLM call, tool invocation, and agent event loop cycle. These are exported to Langfuse~\cite{langfuse2024} via the OTLP protocol, capturing token counts, latencies, model parameters, and full input/output content. All traces within a pipeline run share a session identifier for grouping.

\textbf{Annotation traces} (structured). At each subgraph boundary (scanner, threat generation, attack tree generation, TTP mapping, mitigation generation), we push a Langfuse SDK trace with clean input/output JSON. These traces are tagged with \texttt{annotation} for routing to SME review queues.

%----------------------------------------------------------------------
\subsection{SME Review Pipeline}
%----------------------------------------------------------------------

Annotation traces are routed to Langfuse annotation queues---one per capability. Each queue presents the reviewer with the agent's input context and generated output, along with the applicable scoring dimensions. Reviewers assign categorical scores independently for each dimension and may add free-text feedback.

%----------------------------------------------------------------------
\subsection{TTP Ground Truth Dataset}
%----------------------------------------------------------------------

Individual TTP mappings are automatically pushed as Langfuse dataset items with the attack step description as input and the mapped technique as metadata. The \texttt{expected\_output} field is initially null; SMEs label each mapping as correct or incorrect, building a ground truth corpus over time. This corpus serves dual purposes: (1) evaluating embedding retrieval accuracy via precision@$K$, and (2) providing training pairs for fine-tuning domain-specific embedding models.

%----------------------------------------------------------------------
\subsection{Automated Structural Metrics}
%----------------------------------------------------------------------

We compute the following metrics automatically from pipeline outputs:

\begin{itemize}
    \item \textbf{Tree depth}: Maximum root-to-leaf path length per attack tree.
    \item \textbf{Step count}: Total attack steps across all trees.
    \item \textbf{Technique diversity}: $|\text{unique techniques}| / |\text{total mappings}|$, measuring how broadly the system maps across ATT\&CK.
    \item \textbf{Phase coverage}: Fraction of ATT\&CK tactics (Initial Access through Impact) represented in the generated mappings.
    \item \textbf{Reviewer override rate}: Fraction of TTP mappings where the LLM reviewer changed the embedding's top-1 selection.
\end{itemize}
